\documentclass[10pt,a4paper]{article}
\usepackage[margin=1.9cm]{geometry}
\usepackage{enumitem}
\usepackage[T1]{fontenc}
\usepackage{parskip}
\usepackage[hidelinks]{hyperref}

\setlist[itemize]{topsep=2pt,itemsep=1pt,parsep=0pt,leftmargin=1.4em}
\setlist[enumerate]{topsep=2pt,itemsep=1pt,parsep=0pt,leftmargin=1.6em}
\newcommand{\q}[1]{\medskip\noindent\textbf{#1}\par\nopagebreak}

\pagestyle{empty}

\begin{document}

\begin{center}
{\large\bfseries Information Sheet}\\[2pt]
{\itshape Machine Learning} (Springer)\\[6pt]
\textbf{Corpus-Specific Embedding-Model Selection for Retrieval-Augmented Generation:\\
A Local Benchmarking Pipeline with Query-Level Uncertainty}\\[2pt]
Murat Karakaya, Department of Software Engineering, TED University, Ankara, T\"urkiye\\
\href{mailto:murat.karakaya@tedu.edu.tr}{murat.karakaya@tedu.edu.tr}
\end{center}

\q{1. What is the main claim of the paper? Why is this an important contribution to the
machine learning / data mining research area?}

The main claim is that embedding-model selection for retrieval-augmented generation (RAG)
should be decided per corpus, and that this decision can be made repeatable and statistically
interpretable by an end-to-end local benchmarking pipeline. The pipeline turns a folder of
documents into a complete evaluation benchmark---chunks, synthetic queries, LLM-constructed
relevance judgments, multi-model retrieval runs, uncertainty estimates, and automated
reports---without requiring the practitioner to supply any labeled data.

This matters because standard practice is to pick an embedding model from a public leaderboard.
Public benchmarks are essential shared references, but they measure behavior on public tasks,
not on the practitioner's own documents, and they cannot say whether an observed ranking gap on
a local corpus is large enough to act on. The contribution is a reusable measurement instrument
rather than a new model or metric.

\q{2. What is the evidence provided to support claims? Be precise.}

The pipeline is demonstrated through one frozen experimental snapshot, reported in full and
released as artifacts:

\begin{itemize}
  \item \textbf{Scope.} Three corpora (\texttt{education\_eng}, \texttt{education\_tr},
        \texttt{law\_tr}); 60 source documents; 1{,}914 chunks; five sentence-transformer
        candidates (M1--M5); 150 generated synthetic queries; 9{,}296 judged query--chunk
        pairs, of which 952 are positive; 148 positive-bearing queries scored after
        zero-positive exclusion.
  \item \textbf{Relevance judgments.} Model-independent pooled qrels built from the union of
        top-20 candidates retrieved by all five models, with explicit negative labels retained
        so the judged set is not biased toward any one candidate.
  \item \textbf{Uncertainty.} Percentile and paired bootstrap intervals
        ($B = 10{,}000$) over each corpus's own scored queries. The paired intervals are the
        discriminating evidence: on \texttt{education\_eng} the M1$-$M2 interval includes zero,
        so that corpus does not separate the two leading candidates on retrieval quality and
        the decision falls to latency; on both Turkish corpora every paired interval against
        M5 excludes zero.
  \item \textbf{Pool completeness.} Seed-chunk coverage is measured directly rather than
        assumed: 9 of 150 seed chunks fall outside the pooled candidate set.
  \item \textbf{Operational cost.} Compound corpus embedding and index-insertion throughput,
        and compound query-embedding and vector-search latency, from a single-run CPU-only
        snapshot.
  \item \textbf{Verification.} Every hand-written number in the manuscript is checked against
        the stored run artifacts by an automated gate (33 values); the reported snapshot is
        frozen and released.
\end{itemize}

The paper is explicit about what this evidence does not establish: queries are synthetic and
qrels are LLM-generated without human validation; each language and domain condition is
represented by one corpus, so both are confounded with corpus identity; retrieval geometry,
chunking, pool depth, and cutoff are fixed demonstration settings with no sensitivity study; and
timing is single-run and single-machine. These limits appear in the abstract and Section~6.4.

\q{3. Report the 3--5 most closely related contributions in the past 7 years, and briefly state
the relation of the submission to them.}

\begin{enumerate}
  \item \textbf{Muennighoff et al. (2023), MTEB}, EACL; and \textbf{Enevoldsen et al. (2025),
        MMTEB}, ICLR. These supply corpus, queries, and judgments together and answer which
        model performs best on fixed public tasks. This submission is upstream and
        complementary: it \emph{produces} a labeled benchmark from a practitioner's own
        unlabeled documents---precisely the input MTEB-style evaluation presupposes.
  \item \textbf{Thakur et al. (2025), Assessing Support for the TREC 2024 RAG Track}, SIGIR.
        Characterizes when LLM relevance assessment agrees with human assessment. This
        submission depends on LLM judging and cites this work as the basis for treating its own
        unvalidated qrels as a stated limitation rather than as ground truth.
  \item \textbf{Upadhyay et al. (2024), UMBRELA}. Provides an LLM relevance judge as a
        component. This submission uses a judge of that kind inside a larger corpus-to-report
        workflow and contributes the surrounding pooling, uncertainty, and reporting stages.
  \item \textbf{Es et al. (2024), RAGAs}, EACL Demos. Scores an assembled generation pipeline on
        faithfulness and answer relevance, and can synthesize evaluation data from documents.
        This submission sits upstream: it compares ranked retrieval over the full chunk set, the
        stage RAGAS assumes is already fixed.
  \item \textbf{Wang et al. (2025), On the Reproducibility Limitations of RAG Systems}.
        Documents how RAG evaluation results fail to reproduce. This submission responds by
        freezing and releasing every intermediate artifact and verifying all reported numbers
        against them automatically.
\end{enumerate}

\q{4. Specify 5 general keywords and 5 specific keywords describing the main research activity.}

\textbf{General:} information retrieval; retrieval-augmented generation; evaluation
methodology; benchmarking; reproducibility.

\textbf{Specific:} corpus-specific embedding-model selection; synthetic query generation;
LLM-as-judge relevance construction; model-independent pooled qrels; query-level paired
bootstrap intervals.

\q{5. Conflicts of interest}

\textbf{Institutional conflicts.} The author is affiliated with the Department of Software
Engineering, TED University, Ankara, T\"urkiye. Reviewers from TED University should not be
assigned.

\textbf{Guest Editorial Board conflicts.} None known.

\textbf{Financial and personal conflicts.} None. This research received no specific grant from
funding agencies in the public, commercial, or not-for-profit sectors. The author declares no
competing financial interests or personal relationships that could have appeared to influence
the work.

\textbf{Prior publication.} A preliminary version was presented at the International Conference
on Computer Science and Intelligence Systems (ICCSIS 2025), Madrid, Spain, 10--11 August 2025,
pp.~1--6; the conference organizers subsequently published accepted papers in the
\emph{International Journal of Advances in Electronics and Computer Science} 12(9), ISSN
2394-2835. That version established feasibility on a single English corpus. The present
submission extends it with three language- and domain-differentiated corpora,
model-independent pooled qrels, query-level bootstrap uncertainty, dataset-level reports as the
primary selection artifact, and operational-cost evidence. The relationship is stated in
Section~1 and in full in Section~2.6, and the prior paper is cited and supplied with this
submission. The manuscript is not under review at any other journal or conference.

\q{6. Suggested reviewers (up to four)}

\begin{enumerate}
  \item \textbf{Assoc.\ Prof.\ Dr.\ Beytullah Y{\i}ld{\i}z}, Turkish--Japanese Science and
        Technology University, Ankara, T\"urkiye ---
        \href{mailto:beytullah.yildiz@tju.edu.tr}{beytullah.yildiz@tju.edu.tr}.\\
        Works on machine learning and the design of evaluation metrics for natural-language
        systems, including model-based evaluation of question-answering systems and of
        abstractive summarization, and has supervised graduate work on retrieval-augmented
        generation --- directly relevant to this paper's evaluation methodology.
  \item \textbf{Prof.\ Dr.\ \c{S}eref Sa\u{g}{\i}ro\u{g}lu}, Department of Computer Engineering,
        Gazi University, Ankara, T\"urkiye ---
        \href{mailto:ss@gazi.edu.tr}{ss@gazi.edu.tr}.\\
        Directs the Artificial Intelligence and Big Data Center at Gazi University and works on
        artificial intelligence, big data analytics, and software engineering --- well placed to
        assess the pipeline as an applied machine-learning measurement instrument and its
        reproducibility claims.
  \item \textbf{Prof.\ Dr.\ Murat Koyuncu}, Department of Computer Engineering,
        \c{C}ankaya University, Ankara, T\"urkiye ---
        \href{mailto:murat.koyuncu@cankaya.edu.tr}{murat.koyuncu@cankaya.edu.tr}.\\
        Works on databases, knowledge-based systems, and applied natural language processing
        including Turkish-language sentiment analysis --- relevant to the Turkish corpora and
        to the qrels-construction stage of this paper.
  \item \textbf{Dr.\ Arda Sezen}, At{\i}l{\i}m University, Ankara, T\"urkiye ---
        \href{mailto:arda.sezen@atilim.edu.tr}{arda.sezen@atilim.edu.tr}.\\
        Works on artificial intelligence and deep-learning applications, with published work on
        deep-learning image classification and on the comparative assessment of programming
        environments for artificial-intelligence workloads.
\end{enumerate}

\end{document}
